% ===========================================================================
% Conclusiones
% ===========================================================================

\chapter{Conclusiones}
\label{ch:conclusiones}

La presente investigación tuvo como objetivo mitigar el efecto de los retardos asimétricos en el sincronismo de redes inalámbricas empleando una versión mejorada del protocolo gPTP. A partir del trabajo desarrollado, se establecen las siguientes conclusiones:

\begin{enumerate}
\item \textbf{Revisión del estado del arte.} Se realizó una revisión exhaustiva de las bases teóricas relacionadas con el sincronismo en redes inalámbricas, cubriendo los principales conceptos (offset, skew, tipos de relojes), las aplicaciones en IoT y la Industria 4.0, los desafíos técnicos y los protocolos de sincronización existentes. Se confirmó que gPTP (IEEE 802.1AS) es el más apropiado para redes inalámbricas dentro del ecosistema TSN, aunque su rendimiento se ve significativamente degradado por las asimetrías en los canales de comunicación.

\item \textbf{Análisis sistemático de métodos de corrección.} Se realizó una revisión de 15 métodos publicados entre 2014 y 2026 para mitigar el efecto de los retardos asimétricos en PTP/gPTP, organizados en siete categorías y evaluados en cinco criterios (precisión, complejidad, compatibilidad con gPTP, manejo de asimetrías desconocidas y madurez).

\item \textbf{Selección del método híbrido Exel + AKF.} Se propuso un enfoque híbrido que combina la corrección determinista de Exel (etapa~1) con un AKF de tres estados (etapa~2) para la estimación conjunta de offset ($\theta_k$), skew ($s_k$) y asimetría residual ($\Delta_k$), con adaptación dinámica de la varianza de medición $R_k$.

\item \textbf{Implementación del simulador gPTP.} Se desarrolló un simulador de eventos discretos en MATLAB/GNU Octave, compuesto por siete módulos de código, que modela el comportamiento del estándar IEEE 802.1AS incluyendo el mecanismo \textit{peer-to-peer}, el \textit{rate ratio}, los retardos de procesamiento asimétricos $p_1$--$p_4$, y las imperfecciones de los osciladores.

\item \textbf{Validación experimental completa.} Los resultados de la simulación Monte Carlo ($N = 500$) confirman que la corrección Exel reduce el error en 49.78~$\mu$s (33.2\%, 150.11 $\rightarrow$ 100.33~$\mu$s), reproduciendo fielmente el trabajo de referencia. El método híbrido Exel+AKF propuesto mejora adicionalmente la precisión en un 33.6\% sobre Exel (100.33 $\rightarrow$ 66.24~$\mu$s), alcanzando una reducción total del error del 55.9\% respecto al protocolo sin corrección. La mejora es estadísticamente muy significativa ($t = 25.14$, $p < 0.001$, prueba $t$ pareada), con un \textit{overhead} adicional de apenas el 2\%.

\item \textbf{Contribución metodológica.} La principal contribución es el diseño, la implementación y la validación experimental de una arquitectura de dos etapas (determinista + adaptativa) para la compensación de asimetrías en gPTP que no requiere modificaciones al estándar, es compatible con estampado por software, y se adapta automáticamente a condiciones cambiantes del canal.

\item \textbf{Resultados que superan las proyecciones.} Los resultados validados superan las proyecciones iniciales: la mejora real del AKF (33.6\%) excede la proyectada (28.8\%), mientras que el \textit{overhead} real ($\sim$2\%) es inferior al estimado (5--8\%).

\item \textbf{Limitaciones.} El método asume ruido Gaussiano, aproximación razonable pero no exacta para todas las condiciones de canal inalámbrico. El estudio se limita a un enlace punto a punto de un solo salto, y los parámetros del AKF pueden requerir reajuste para escenarios con características de canal significativamente diferentes.
\end{enumerate}
